\hmsubsection{Embedded Systems}{UMA}{OAH}

An embedded system is a computing system designed to perform a 
dedicated function within a larger mechanical or electrical system, 
typically operating under constraints related to processing power, 
memory, energy consumption, and real-time response \cite{marwedel2021}. 
Unlike general-purpose computers, embedded systems are tightly coupled 
to the physical environment in which they operate, and their correctness 
depends not only on the logical result of computation but also on the 
time at which the result is produced \cite{lee2017}.

In the context of this project, the robotic system follows the 
sense--process--act cycle described in Section~\ref{sec:robotics}, and each stage is 
implemented on dedicated computing hardware. The camera acquires image 
data, a host processor analyses the data and computes motor positions, 
and a microcontroller translates these positions into electrical 
commands that drive the actuators. Because the actuators are physical 
motors that must respond predictably and within bounded time, the 
control path is subject to real-time constraints that cannot be 
guaranteed by a general-purpose operating system alone.

To address this, the system is implemented using a two-tier 
architecture consisting of a higher-level host processor and a 
lower-level microcontroller. The host processor is a Raspberry Pi~5 
running a Linux-based operating system, on which the image processing 
pipeline and the inverse kinematics solver are executed. The 
microcontroller is an OpenRB-150 board running bare-metal C++ firmware, 
dedicated exclusively to driving the five Dynamixel servos over a TTL 
bus. Communication between the two tiers is performed over USB serial.

This separation of concerns is a common pattern in embedded robotic 
systems \cite{lee2017}. The host processor provides the computational 
resources required for vision and kinematics, including access to 
standard libraries such as OpenCV and NumPy, while the microcontroller 
provides the deterministic timing required for stable motor control. 
By isolating the time-critical motor control loop on dedicated hardware, 
the system avoids the timing variability introduced by the host's 
process scheduler, memory management, and other concurrent tasks.

The hardware components, communication protocol, and firmware that 
realise this architecture are described in detail in 
Sections~\ref{sec:hardware-architecture}, \ref{sec:comm}, 
and~\ref{sec:actuator-control}.
